% Chapter converted from Markdown
% Auto-generated - do not edit manually

\section*{Chapter 21 - Confidence Calibration}


\subsection*{Purpose}
\item Detail the mathematical calibration of confidence signals (VIF, authority, capability) to keep decisions reliable.
\item Describe Bayesian update routines, calibration experiments, and dashboards that track confidence accuracy.
\item Provide runnable commands to observe calibration data and adjust confidence.

\subsection*{Calibration Model}

Confidence is treated as a probability distribution updated via Bayes' rule:
``\texttt{
posterior = prior * likelihood / evidence
}`\texttt{

\textbf{Component Details:}

\textbf{Prior Distribution:}
\item Historical accuracy of the system/persona in similar contexts
\item Formula: }prior = Beta(α, β)\texttt{ where α = successes, β = failures
\item Updated after each outcome: }α\_new = α + success\texttt{, }β\_new = β + failure\texttt{
\item Mean: }E[prior] = α / (α + β)\texttt{

\textbf{Likelihood Function:}
\item Probability of observing current evidence if claim is correct
\item Sources: quality gates (SDF-CVF), tests (quartet parity), audits (CAS)
\item Formula: }likelihood = P(evidence | claim\_true)\texttt{
\item Combined: }likelihood = ∏ P(gate\_i | claim\_true)\texttt{ for all gates i

\textbf{Evidence (Normalization):}
\item Ensures probabilities sum to 1.0
\item Formula: }evidence = prior \textit{ likelihood + (1 - prior) } (1 - likelihood)\texttt{
\item Marginal probability of observing the evidence

\textbf{Posterior Distribution:}
\item Updated confidence after observing evidence
\item Formula: }posterior = (prior * likelihood) / evidence\texttt{
\item Mean: }E[posterior] = (α + successes) / (α + β + total\_observations)\texttt{

\textbf{Calibration Curves:}
\item Map predicted confidence to observed success rates
\item Bins: [0.0-0.1], [0.1-0.2], ..., [0.9-1.0]
\item For each bin: }calibration = observed\_success\_rate / predicted\_confidence\texttt{
\item Perfect calibration: calibration = 1.0 for all bins
\item Deviations trigger rebalancing via prior updates

\subsection*{Confidence Types}

AIM-OS distinguishes four confidence types to prevent inflation:

\textbf{Type 1: Direction Confidence}
\item Question: "Is this the RIGHT choice?"
\item Example: VIF implementation = 0.95 (clearly serves OBJ-03)
\item Context: Strategic alignment

\textbf{Type 2: Execution Confidence}
\item Question: "Can I DO this successfully?"
\item Example: VIF implementation = 0.65 (never built code from docs)
\item Context: Technical capability

\textbf{Type 3: Autonomous Confidence}
\item Question: "Can I do this ALONE without help?"
\item Example: VIF implementation = 0.60 (will need questions answered)
\item Context: Self-sufficiency

\textbf{Type 4: Collaborative Confidence}
\item Question: "Can I do this WITH support?"
\item Example: VIF implementation = 0.75 (can ask when stuck)
\item Context: Team collaboration

\textbf{Calibration Model Per Type:}
}`\texttt{python
calibration\_model = {
    "documentation\_tasks": {
        "bias": -0.05,  \# Slightly underconfident
        "accuracy": 0.95  \# Usually correct
    },
    "code\_tasks": {
        "bias": +0.20,  \# OVERCONFIDENT (predicted 0.85, actual 0.65)
        "accuracy": 0.70  \# Sometimes struggle
    },
    "organizational\_tasks": {
        "bias": 0.00,  \# Well calibrated
        "accuracy": 0.90
    }
}

\section*{When making new decision:}
raw\_confidence = my\_intuition()  \# 0.85
task\_category = classify(task)  \# "code\_tasks"
calibrated\_confidence = raw\_confidence - calibration\_model[task\_category]["bias"]
\section*{0.85 - 0.20 = 0.65 ← HONEST confidence}
}`\texttt{

\subsection*{Runnable Examples}

\subsubsection*{Example 1: Read Calibration Metrics}
}`\texttt{powershell
\section*{Read current confidence metrics (includes calibration summary)}
\_\_MATH\_BLOCK\_0\_\_true;
        include\_ece=\_\_MATH\_BLOCK\_1\_\_result = Invoke-WebRequest -Uri 'http://localhost:5001/mcp/execute' }
    -Method POST -ContentType 'application/json' -Body \_\_MATH\_BLOCK\_2\_\_(\_\_MATH\_BLOCK\_3\_\_(\_\_MATH\_BLOCK\_4\_\_result.calibration.bins | ForEach-Object {
    Write-Host "  [\_\_MATH\_BLOCK\_5\_\_\_.range)]: Predicted=\_\_MATH\_BLOCK\_6\_\_\_.predicted), Observed=\_\_MATH\_BLOCK\_7\_\_\_.observed), Count=\_\_MATH\_BLOCK\_8\_\_\_.count)"
}
``\texttt{

\subsubsection*{Example 2: Track Confidence Update}
}`\texttt{powershell
\section*{Record a confidence update after validation}
\_\_MATH\_BLOCK\_9\_\_update |
    Select-Object -ExpandProperty Content | ConvertFrom-Json
``\texttt{

\subsubsection*{Example 3: Calculate Calibration Bias}
}`\texttt{powershell
\section*{Query confidence history for bias calculation}
\_\_MATH\_BLOCK\_10\_\_messages = Invoke-WebRequest -Uri 'http://localhost:5001/mcp/execute' }
    -Method POST -ContentType 'application/json' -Body \_\_MATH\_BLOCK\_11\_\_bias = @{}
\_\_MATH\_BLOCK\_12\_\_\_.tags.task\_category } | Group-Object { \_\_MATH\_BLOCK\_13\_\_category = \_\_MATH\_BLOCK\_14\_\_predictions = \_\_MATH\_BLOCK\_15\_\_\_.tags.predicted\_confidence }
    \_\_MATH\_BLOCK\_16\_\_\_.Group | ForEach-Object { \_\_MATH\_BLOCK\_17\_\_bias[\_\_MATH\_BLOCK\_18\_\_predictions | Measure-Object -Average).Average - (\_\_MATH\_BLOCK\_19\_\_bias.GetEnumerator() | ForEach-Object {
    Write-Host "  \_\_MATH\_BLOCK\_20\_\_\_.Key): \_\_MATH\_BLOCK\_21\_\_\_.Value)"
}
``\texttt{

\subsection*{Calibration Workflow}
1. Record prediction (confidence before execution).
2. Run validations (examples, audits, deployments).
3. Compare outcome with prediction; compute calibration error.
4. Update priors and weightings; log results in CAS + SEG.

\subsection*{Expected Calibration Error (ECE)}

ECE measures how well-calibrated probabilistic predictions are:

\textbf{Formula:}
}`\texttt{
ECE = Σ (|B\_m| / n) * |acc(B\_m) - conf(B\_m)|
}`\texttt{

Where:
\item }B\_m\texttt{: Bin m containing predictions
\item }|B\_m|\texttt{: Number of predictions in bin m
\item }n\texttt{: Total number of predictions
\item }acc(B\_m)\texttt{: Actual accuracy in bin m
\item }conf(B\_m)\texttt{: Average predicted confidence in bin m

\textbf{Interpretation:}
\item ECE = 0.0: Perfect calibration (predicted = actual)
\item ECE < 0.05: Well calibrated
\item ECE > 0.10: Poor calibration (requires adjustment)

\textbf{Binning Strategy:}
\item Optimal bin count: }sqrt(n)\texttt{ where n = number of predictions
\item Equal-width bins: [0.0-0.1], [0.1-0.2], ..., [0.9-1.0]
\item Equal-frequency bins: Each bin contains equal number of predictions

\subsection*{Metrics \& Dashboards}

\textbf{Brier Score:}
\item Measures accuracy of probabilistic predictions
\item Formula: }Brier = (1/n) * Σ (predicted\_i - actual\_i)²\texttt{
\item Range: [0.0, 1.0] where 0.0 = perfect, 1.0 = worst
\item Decomposes into: }Brier = Calibration + Resolution + Uncertainty\texttt{
  - Calibration: How well probabilities match frequencies
  - Resolution: How well predictions distinguish outcomes
  - Uncertainty: Inherent unpredictability

\textbf{Calibration Bins:}
\item Bucket predictions into bins (e.g., 0.5-0.6, 0.6-0.7, etc.)
\item Compare expected vs actual success rates
\item Visualize as calibration plot: predicted (x-axis) vs observed (y-axis)
\item Perfect calibration: diagonal line (y = x)

\textbf{VIF Drift:}
\item Monitors changes after calibration updates
\item Formula: }drift = |VIF\_new - VIF\_old|\texttt{
\item Threshold: drift > 0.10 triggers review
\item Tracks: per-system, per-persona, per-task-type

\textbf{Confidence Gap Log:}
\item Highlights systems consistently over/under confident
\item Overconfidence: predicted > actual (bias > 0)
\item Underconfidence: predicted < actual (bias < 0)
\item Tracks: bias per task category, temporal trends, improvement velocity

\subsection*{Failure Modes \& Mitigations}
\item \textbf{Overconfidence:} tighten gates; require additional evidence; adjust priors.
\item \textbf{Underconfidence:} encourage more automation; add proof tasks; revisit penalties.
\item \textbf{Data sparsity:} increase sampling; run synthetic experiments; aggregate across similar contexts.
\item \textbf{Model drift:} retrain calibration curves; run regression tests.

\subsection*{Integration}
\item \textbf{VIF:} uses calibrated confidence to gate work; stores updates per chapter/system.
\item \textbf{CAS:} awareness dashboards display calibration status; anomalies trigger alerts.
\item \textbf{SDF-CVF:} quality results feed likelihood calculations.
\item \textbf{APOE/SIS:} improvements proposed when calibration error exceeds thresholds.

\subsection*{Mathematical Foundations}

Confidence calibration in AIM-OS is grounded in Bayesian probability theory and statistical learning. The mathematical framework ensures that confidence signals accurately reflect the true probability of success.

\subsubsection*{Bayesian Update Theory}

\textbf{Bayesian Inference Framework:}
\item \textbf{Prior Belief:} Initial confidence based on historical performance
\item \textbf{Evidence:} Observed outcomes from quality gates, tests, audits
\item \textbf{Posterior Belief:} Updated confidence after observing evidence
\item \textbf{Conjugate Prior:} Beta distribution for binary outcomes (success/failure)

\textbf{Beta Distribution Properties:}
\item Parameters: α (successes), β (failures)
\item Mean: }E[θ] = α / (α + β)\texttt{
\item Variance: }Var[θ] = (α × β) / ((α + β)² × (α + β + 1))\texttt{
\item Mode: }(α - 1) / (α + β - 2)\texttt{ for α, β > 1
\item Conjugate: Beta-Binomial conjugacy enables efficient updates

\textbf{Update Rule:}
}`\texttt{
α\_new = α\_old + successes
β\_new = β\_old + failures
}`\texttt{

\textbf{Confidence Interval:}
\item 95\% credible interval: }[Beta(α, β).ppf(0.025), Beta(α, β).ppf(0.975)]\texttt{
\item Provides uncertainty quantification alongside point estimate

\subsubsection*{Calibration Theory}

\textbf{Perfect Calibration Definition:}
A system is perfectly calibrated if, for all confidence levels c:
}`\texttt{
P(success | predicted\_confidence = c) = c
}`\texttt{

\textbf{Calibration Error:}
\item Measures deviation from perfect calibration
\item Formula: }CE = E[|predicted - actual|]\texttt{
\item Decomposes into: systematic bias + random error

\textbf{Calibration Curve:}
\item Maps predicted confidence to observed success rate
\item Perfect calibration: diagonal line (y = x)
\item Overconfidence: curve below diagonal
\item Underconfidence: curve above diagonal

\subsubsection*{Statistical Learning Framework}

\textbf{Online Learning:}
\item Updates calibration model after each outcome
\item Exponential moving average: }θ\_t = α × θ\_{t-1} + (1 - α) × x\_t\texttt{
\item Adaptive learning rate: decreases over time for stability

\textbf{Task Category Clustering:}
\item Groups similar tasks for bias estimation
\item Features: task type, complexity, domain, tools used
\item Clustering: k-means or hierarchical clustering
\item Bias per cluster: }bias\_k = mean(predicted\_k - actual\_k)\texttt{

\subsection*{Calibration Algorithms}

AIM-OS implements several calibration algorithms to maintain accurate confidence signals:

\subsubsection*{Algorithm 1: Beta-Binomial Calibration}

\textbf{Purpose:} Update confidence using Beta-Binomial conjugacy

\textbf{Algorithm:}
}`\texttt{python
def beta\_binomial\_calibration(prior\_alpha, prior\_beta, successes, failures):
    """
    Update Beta prior with observed outcomes.
    
    Args:
        prior\_alpha: Prior α parameter (successes)
        prior\_beta: Prior β parameter (failures)
        successes: Observed successes
        failures: Observed failures
    
    Returns:
        Updated (α, β) parameters
    """
    alpha\_new = prior\_alpha + successes
    beta\_new = prior\_beta + failures
    return alpha\_new, beta\_new

\section*{Example usage}
alpha, beta = beta\_binomial\_calibration(α=10, β=2, successes=5, failures=1)
confidence = alpha / (alpha + beta)  \# Updated confidence
}`\texttt{

\textbf{Properties:}
\item Efficient: O(1) update time
\item Memory: Stores only (α, β) parameters
\item Interpretable: Direct probability interpretation

\subsubsection*{Algorithm 2: Platt Scaling}

\textbf{Purpose:} Calibrate confidence scores using logistic regression

\textbf{Algorithm:}
}`\texttt{python
def platt\_scaling(raw\_scores, labels):
    """
    Calibrate raw confidence scores using Platt scaling.
    
    Args:
        raw\_scores: Raw confidence scores [0, 1]
        labels: Binary outcomes (0 = failure, 1 = success)
    
    Returns:
        Calibrated scores
    """
    from sklearn.calibration import CalibratedClassifierCV
    from sklearn.linear\_model import LogisticRegression
    
    \# Fit Platt scaling
    platt = CalibratedClassifierCV(
        LogisticRegression(),
        method='sigmoid',
        cv=5
    )
    platt.fit(raw\_scores.reshape(-1, 1), labels)
    
    \# Calibrate new scores
    calibrated = platt.predict\_proba(raw\_scores.reshape(-1, 1))[:, 1]
    return calibrated
}`\texttt{

\textbf{Properties:}
\item Non-parametric: No distribution assumptions
\item Flexible: Adapts to any score distribution
\item Requires: Calibration dataset

\subsubsection*{Algorithm 3: Isotonic Regression}

\textbf{Purpose:} Non-parametric calibration using isotonic regression

\textbf{Algorithm:}
}`\texttt{python
def isotonic\_calibration(raw\_scores, labels):
    """
    Calibrate using isotonic regression (piecewise constant).
    
    Args:
        raw\_scores: Raw confidence scores [0, 1]
        labels: Binary outcomes (0 = failure, 1 = success)
    
    Returns:
        Calibrated scores
    """
    from sklearn.isotonic import IsotonicRegression
    
    \# Fit isotonic regression
    iso = IsotonicRegression(out\_of\_bounds='clip')
    iso.fit(raw\_scores, labels)
    
    \# Calibrate new scores
    calibrated = iso.transform(raw\_scores)
    return calibrated
}`\texttt{

\textbf{Properties:}
\item Non-parametric: No distribution assumptions
\item Monotonic: Preserves score ordering
\item Flexible: Piecewise constant mapping

\subsubsection*{Algorithm 4: Temperature Scaling}

\textbf{Purpose:} Single-parameter calibration for neural network outputs

\textbf{Algorithm:}
}`\texttt{python
def temperature\_scaling(logits, temperature):
    """
    Calibrate logits using temperature scaling.
    
    Args:
        logits: Raw logits from model
        temperature: Temperature parameter (T > 0)
    
    Returns:
        Calibrated probabilities
    """
    import torch.nn.functional as F
    
    \# Apply temperature scaling
    scaled\_logits = logits / temperature
    calibrated = F.softmax(scaled\_logits, dim=-1)
    return calibrated

\section*{Optimal temperature via cross-validation}
def find\_optimal\_temperature(logits, labels, temp\_range=[0.1, 10.0]):
    """
    Find optimal temperature via cross-validation.
    """
    best\_temp = 1.0
    best\_ece = float('inf')
    
    for temp in np.linspace(temp\_range[0], temp\_range[1], 100):
        calibrated = temperature\_scaling(logits, temp)
        ece = calculate\_ece(calibrated, labels)
        if ece < best\_ece:
            best\_ece = ece
            best\_temp = temp
    
    return best\_temp
}`\texttt{

\textbf{Properties:}
\item Simple: Single parameter to tune
\item Efficient: O(n) computation
\item Effective: Works well for neural networks

\subsection*{System Architecture}

The confidence calibration system integrates with all AIM-OS systems to provide accurate confidence signals:

\subsubsection*{Core Components}

\textbf{1. Calibration Tracker}
\item \textbf{Purpose:} Records predicted confidence and actual outcomes
\item \textbf{Storage:} CMC atoms tagged }confidence\_calibration\texttt{
\item \textbf{Schema:} }{task\_id, predicted, actual, task\_category, timestamp, confidence\_type}\texttt{
\item \textbf{Updates:} Real-time updates after each task completion

\textbf{2. Bias Calculator}
\item \textbf{Purpose:} Calculates calibration bias per task category
\item \textbf{Algorithm:} Mean difference between predicted and actual
\item \textbf{Output:} Bias per category: }bias\_k = mean(predicted\_k - actual\_k)\texttt{
\item \textbf{Updates:} Recalculated after each batch of outcomes

\textbf{3. Calibration Model}
\item \textbf{Purpose:} Stores calibration parameters per task category
\item \textbf{Storage:} CMC atoms tagged }calibration\_model\texttt{
\item \textbf{Schema:} }{category, bias, accuracy, sample\_size, last\_updated}\texttt{
\item \textbf{Usage:} Applied to raw confidence before reporting

\textbf{4. ECE Calculator}
\item \textbf{Purpose:} Computes Expected Calibration Error
\item \textbf{Algorithm:} Binned calibration error calculation
\item \textbf{Binning:} Optimal bin count: }sqrt(n)\texttt{ where n = predictions
\item \textbf{Output:} ECE score and calibration curve

\textbf{5. Dashboard Generator}
\item \textbf{Purpose:} Generates calibration dashboards and reports
\item \textbf{Visualizations:} Calibration curves, bias plots, ECE trends
\item \textbf{Alerts:} Triggers when ECE > 0.10 or bias exceeds thresholds
\item \textbf{Integration:} CAS dashboards, VIF confidence displays

\subsubsection*{Data Flow}

\textbf{Calibration Flow:}
}`\texttt{
Task Start → Record Predicted Confidence → Execute Task → 
Record Actual Outcome → Calculate Error → Update Calibration Model → 
Apply Calibration to Future Predictions
}`\texttt{

\textbf{Bias Calculation Flow:}
}`\texttt{
Collect Outcomes → Group by Task Category → Calculate Mean Error → 
Update Bias Model → Store in CMC → Apply to Raw Confidence
}`\texttt{

\textbf{ECE Calculation Flow:}
}`\texttt{
Collect Predictions → Bin by Confidence Level → Calculate Observed Rate → 
Compare to Predicted → Compute ECE → Generate Calibration Curve
}``

\subsection*{Operational Guidance}

\subsubsection*{Calibration Best Practices}

\textbf{1. Regular Calibration Updates}
\item Update calibration models after every 20-50 outcomes
\item Recalculate bias per category weekly
\item Recompute ECE monthly
\item Review calibration curves quarterly

\textbf{2. Task Category Management}
\item Define clear task categories (documentation, code, organizational)
\item Ensure sufficient samples per category (minimum 10 outcomes)
\item Merge similar categories if sample size too small
\item Split categories if bias varies significantly within category

\textbf{3. Confidence Type Selection}
\item Use Direction Confidence for strategic decisions
\item Use Execution Confidence for technical tasks
\item Use Autonomous Confidence for independent work
\item Use Collaborative Confidence for team tasks

\textbf{4. Calibration Thresholds}
\item ECE < 0.05: Well calibrated, continue monitoring
\item ECE 0.05-0.10: Acceptable, minor adjustments needed
\item ECE > 0.10: Poor calibration, requires recalibration
\item Bias > 0.15: Significant bias, immediate correction needed

\subsubsection*{Troubleshooting Guide}

\textbf{Problem: High ECE (> 0.10)}
\item \textbf{Causes:} Model drift, insufficient data, category mismatch
\item \textbf{Solutions:} Recalibrate model, increase sample size, refine categories

\textbf{Problem: Persistent Overconfidence}
\item \textbf{Causes:} Optimistic priors, insufficient penalty for failures
\item \textbf{Solutions:} Adjust priors downward, increase failure weight, tighten gates

\textbf{Problem: Persistent Underconfidence}
\item \textbf{Causes:} Pessimistic priors, excessive penalty for failures
\item \textbf{Solutions:} Adjust priors upward, decrease failure weight, encourage automation

\textbf{Problem: Category-Specific Bias}
\item \textbf{Causes:} Different difficulty levels, different success criteria
\item \textbf{Solutions:} Split categories, adjust category-specific priors, refine task classification

\subsection*{Mathematical Foundations}
